\documentclass[10pt]{article}
\usepackage{icmlko}

\icmlmeta{IP 파운데이션 월드모델}{261}{실린 무게는 거짓말한다}{하마터면 맞는 노트를 틀렸다고 고칠 뻔했다 --- SVD 무게가 아니라 떼어 보기가 답한다}

\begin{document}

\twocolumn[
\icmlheader

\begin{abstract}
노트 260이 웹툰 메타 축으로 판을 0.4687 $\to$ 0.4916 으로 올리고 그것을
``\textbf{어느 요일인가}''라고 불렀다. 그 절차를 애니에 옮겨 봤다 --- 먼저
사후 후보를 훑어 \texttt{is\_ending} 을 걸러내고(나머지 여덟은 완결과
$|\rho| \le 0.07$ 로 깨끗했다) 성분 검사를 걸었다. \textbf{안 된다}: 통과
성분이 하나 나오지만 판이 $+$0.0029 이고 \textbf{진짜 총량이 0.0000} ---
$|t| \ge 2$ 인 도메인이 하나도 없다. 노트 245의 선 아래다. 까닭은 겹침 표에
있다: 애니 메타 성분들이 \texttt{target\_breadth} 와 0.766 · \texttt{media\_push}
와 0.727 로 겹친다. \textbf{애니의 \texttt{target\_breadth} 가 바로 연령
등급}이라(\texttt{scale01(age, 0, 19)}) 메타의 알맹이가 이미 판에 있다.
\textbf{그 다음이 이 노트의 본론이다.} 웹툰 메타 성분에 무엇이 실렸나 보려고
SVD 무게를 열었더니 --- 성분 0 은 연령 등급($+$0.568)과 작가 수($-$0.471)가
잡고 있고 \textbf{요일은 1\%}, 성분 3 도 작가 수가 앞이고 요일은 24\% 였다.
``노트 260의 제목이 틀렸다''고 쓸 참이었다. \textbf{쓰기 전에 떼어 봤다.}
요일만 넣으면 판 \textbf{0.4881}(웹툰 0.4603, $t{=}5.2$), 등급만 0.4701,
작가 수만 0.4714, \textbf{등급$+$작가 수는 0.4691 로 아무것도 아니다.}
요일이 거의 다 낸다. 무게가 1\% 인 것이 이득의 85\%를 낸다. 까닭은
단순하다 --- 요일 토큰 일곱은 서로 거의 배타적이라 분산이 여러 성분에
흩어지고, 등급 · 작가 수는 소수 토큰에 뭉쳐 앞 성분을 차지한다.
\textbf{실린 무게는 예측 기여가 아니다.} 노트 260은 그대로 서고, 하마터면
맞는 노트를 틀리게 고칠 뻔했다.
\end{abstract}
]

\section{애니에는 안 통한다}

노트 260의 메타 절차를 애니로 옮겼다. 두 번 데인 뒤라(노트 255의 완결 태그 ·
노트 260의 \texttt{daily\_pass}) \textbf{먼저 사후 검사를 돌렸다.}

\begin{center}\small
\begin{tabular}{lrr}
\hline
후보 & \texttt{is\_ending} 과 & 연도와 \\
\hline
\texttt{age} & $+$0.006 & $-$0.063 \\
\texttt{medium} & $-$0.040 & $+$0.132 \\
\texttt{is\_adult} & $+$0.014 & $-$0.098 \\
\texttt{is\_dubbed} & $+$0.020 & $+$0.014 \\
\texttt{is\_free} & $+$0.016 & $-$0.173 \\
\texttt{is\_only} & $+$0.007 & $+$0.090 \\
\texttt{is\_original} & $-$0.015 & $+$0.056 \\
\textbf{\texttt{is\_ending}} & $+$1.000 & $-$0.297 \\
\hline
\end{tabular}
\end{center}

완결 표지 자신 말고는 깨끗하다. 그런데 결과가 안 나온다.

\begin{center}\small
\begin{tabular}{lrrr}
\hline
 & 판 & 애니 & 진짜 총량 \\
\hline
더빙 넣음 & $+$0.0029 & $+$0.0137 ($t{=}1.4$) & \textbf{0.0000} \\
더빙 뺌 & $+$0.0021 & $+$0.0037 ($t{=}0.3$) & \textbf{0.0000} \\
\hline
\end{tabular}
\end{center}

\textbf{$|t| \ge 2$ 인 도메인이 하나도 없다.} 노트 245의 규약대로 채택 안
한다. 겹침 표가 까닭을 말한다 --- 성분 1 이 \texttt{media\_push} 와 0.727,
성분 2 가 \texttt{target\_breadth} 와 0.766 이다. 그리고 \textbf{애니의
\texttt{target\_breadth} 는 \texttt{scale01(age, 0, 19)}, 곧 연령 등급}이다.
메타의 알맹이가 이미 판에 있었다. 노트 256이 세계애니에서 만난 것과 같은
구조다.

\textbf{메타 요령은 보편이 아니다.} 도메인마다 손 축 다섯이 무엇으로 채워져
있느냐에 달렸다.

\section{하마터면 맞는 노트를 고칠 뻔했다}

그러면 웹툰에서는 무엇이 달랐나. 웹툰 메타 성분에 실린 무게를 열었다.

\begin{center}\small
\begin{tabular}{lrr}
\hline
토큰 & 성분 0 & 성분 3 \\
\hline
\texttt{a:RATE\_15} (등급) & $+$0.568 & --- \\
\texttt{a:ALL} (등급) & $-$0.549 & --- \\
\texttt{n:1} (작가 수) & $-$0.471 & $+$0.237 \\
\texttt{n:3} (작가 수) & --- & $-$0.687 \\
\texttt{d:MONDAY} (요일) & $-$0.072 & $+$0.450 \\
\hline
\textbf{요일이 차지하는 몫} & \textbf{1\%} & \textbf{24\%} \\
\hline
\end{tabular}
\end{center}

노트 260의 제목이 ``어느 요일인가''인데 요일은 성분 0 의 1\% 다. \textbf{제목이
틀렸다고 정정 노트를 쓸 참이었다.} 노트 253에서 안 재고 쓴 문장 하나로
노트를 틀리게 만든 적이 있어서, 이번엔 쓰기 전에 쟀다.

\section{떼어 보니 반대였다}

\begin{figure}[t]
\centering
\includegraphics[width=\columnwidth]{figs/load.pdf}
\caption{왼쪽 --- 토큰 무리를 떼어 넣어 본 판. 오른쪽 --- SVD 무게와 실제
기여가 어긋난다.}
\end{figure}

\begin{center}\small
\begin{tabular}{lrrr}
\hline
넣은 토큰 & 판 & 웹툰 & 웹툰 $t$ \\
\hline
\textbf{셋 다} & \textbf{0.4916} & 0.4768 & 5.33 \\
\textbf{요일만} & \textbf{0.4881} & 0.4603 & \textbf{5.20} \\
요일$+$등급 & 0.4872 & 0.4606 & 5.05 \\
작가 수만 & 0.4714 & 0.4060 & 2.99 \\
등급만 & 0.4701 & 0.4010 & 1.24 \\
\textbf{등급$+$작가 수} & \textbf{0.4691} & 0.3994 & \textbf{0.74} \\
바탕(21축) & 0.4687 & 0.3943 & --- \\
\hline
\end{tabular}
\end{center}

\textbf{요일만으로 $+$0.0194 --- 전체 이득 $+$0.0229 의 85\% 다.} 성분 0 을
통째로 잡고 있던 등급과 작가 수는 \emph{둘이 같이} 넣어도 $+$0.0004 로
아무것도 아니다. \textbf{노트 260의 제목이 맞았다.}

까닭은 SVD 의 성질이다. 요일 토큰 일곱은 한 작품이 한둘만 갖는 \textbf{거의
배타적인} 집합이라 분산이 여러 성분에 얇게 흩어진다. 등급(3종)과 작가
수(3종)는 소수 토큰에 뭉쳐 있어 앞 성분을 통째로 차지한다. \textbf{앞
성분을 차지하는 것과 목표를 설명하는 것은 다른 일이다} --- SVD 는 목표를
한 번도 안 보고 분산만 보기 때문이다(그것이 이 축의 장점이기도 하다).

\section{규약}

\textbf{실린 무게로 기여를 읽지 않는다. 떼어 보고 읽는다.} 이 실험실에는
같은 모양의 규약이 이미 둘 있다 --- 노트 253의 ``자료가 모자라다고 쓰기
전에 덜어 내 본다''와 노트 259의 ``자리가 모자라다고 쓰기 전에 잡음을 먹여
본다''. 셋 다 \textbf{설명을 쓰기 전에 한 줄 재라}는 같은 말이고, 세 번 다
짐작이 틀렸다.

이번 것은 방향이 하나 더 있다. 앞의 둘은 \emph{내 짐작을 지웠고}, 이번 것은
\textbf{내 정정을 지웠다.} 틀린 정정도 틀린 주장만큼 나쁘다.

\paragraph{남은 것}
\begin{center}\small
\begin{tabular}{ll}
\hline
갈래 & 왜 \\
\hline
요일이 왜 세나 & 슬롯 경쟁인가 요일별 트래픽인가 \\
23축 천장 & 또 올랐을 것이다 \\
모바일 · 세계애니 메타 & 애니처럼 이미 판에 있나 \\
성분 대신 토큰 & 요일 일곱을 그냥 일곱 축으로 주면? \\
\hline
\end{tabular}
\end{center}

\paragraph{규약} SVD 성분을 쓸 때 \textbf{무엇이 실렸는지로 이름 붙이지
않는다.} 이름은 떼어 보기로 정한다.

\end{document}
